\documentclass[a4paper,12pt]{article}

\usepackage[czech]{babel}
\usepackage[utf8]{inputenc}
\usepackage[T1]{fontenc}
\usepackage{geometry}
\usepackage{enumitem}

\geometry{margin=2.5cm}

\title{\textbf{Kolej CUP}\\
Pravidla turnaje v pozemáčku}
\date{}

\begin{document}

\maketitle

\section*{1. Obecná ustanovení}
Kolej CUP je přátelský sportovní turnaj v pozemáčku zaměřený na fair play, respekt a zábavu.
Všichni účastníci jsou povinni řídit se těmito pravidly a pokyny organizátorů a rozhodčích.

\section*{2. Formát turnaje}
\begin{itemize}
	\item Turnaje se účastní 10 týmů rozdělených do 2 skupin po 5 týmech.
	\item Ve skupině hraje každý tým s každým.
	\item Po základní části postupuje \textbf{8 týmů do play-off}.
	\item Play-off se hraje vyřazovacím systémem na \textbf{1 vítězný zápas}.
	\item Součástí je i \textbf{zápas o 3. místo}.
\end{itemize}

\section*{3. Herní systém zápasu}
\begin{itemize}
	\item Zápas se hraje na \textbf{2x 10 minut hrubého času}.
	\item Čas se nezastavuje.
	\item Hraje se ve formátu \textbf{3 na 3} bez brankáře.
	\item Maximální počet hráčů v týmu: \textbf{8}.
	\item Střídání probíhá libovolně během hry.
\end{itemize}

\section*{4. Bodování a pořadí ve skupině}
\begin{itemize}
	\item Výhra: \textbf{2 body}
	\item Remíza: \textbf{1 bod}
	\item Prohra: \textbf{0 bodů}
\end{itemize}

Pořadí ve skupině se určuje podle:
\begin{enumerate}
	\item vyššího počtu bodů,
	\item lepších vzájemných zápasů,
	\item lepšího rozdílu skóre,
	\item vyššího počtu vstřelených gólů.
\end{enumerate}

\section*{5. Play-off}
\begin{itemize}
	\item V případě remízy následuje \textbf{prodloužení 222 sekund}.
	\item Pokud nerozhodne prodloužení, následuje \textbf{trestné střílení}.
\end{itemize}

\subsection*{Trestné střílení}
\begin{itemize}
	\item Každý tým nominuje \textbf{3 střelce}.
	\item Pokud není rozhodnuto, pokračuje se systémem náhlé smrti.
\end{itemize}

\section*{6. Herní pravidla}
\begin{itemize}
	\item Hraje se s \textbf{oranžádou}.
	\item Hra je bez brankáře.
	\item Důraz je kladen na plynulost hry a bezpečnost hráčů.
\end{itemize}

\subsection*{6.1 Brankoviště}
\begin{itemize}
	\item Brankoviště je prostor vyznačený křídou před brankou.
	\item ..........Doporučený tvar: půlkruh o poloměru cca \textbf{1,5 m} od středu branky...........
	\item V brankovišti \textbf{nesmí stát ani setrvávat bránící hráč}.
	\item Pokud bránící hráč ovlivní hru z prostoru brankoviště, rozhodčí může:
	      \begin{itemize}
		      \item uznat vstřelený gól,
		      \item nebo nařídit samostatný nájezd.
	      \end{itemize}
	\item Útočící hráč může brankovištěm pouze projíždět, nesmí v něm zůstávat a blokovat hru.
\end{itemize}

\subsection*{6.2 Rozehrávka}
\begin{itemize}
	\item Pokud je míček vyhozen mimo hrací plochu, hra se přerušuje.
	\item Rozehrává \textbf{tým, který se \textbf{ne}dotkl míčku jako poslední}.
	\item Rozehrávka probíhá z místa, kde míček opustil hrací plochu (nebo nejbližšího vhodného bodu).
	\item Hráči soupeře musí stát alespoň \textbf{2 metry} od míčku.
	\item Rozehrávka musí být provedena bez zbytečného zdržování.
	\item Při úmyslném zdržování nebo opakovaném vyhazování může rozhodčí nařídit přísnější trest (např. nájezd).
\end{itemize}

\section*{7. Fauly a tresty}
Za faul se považuje zejména:
\begin{itemize}
	\item hra vysokou holí,
	\item sekání,
	\item nebezpečný zákrok,
	\item kosček,
	\item nedovolené bránění,
	\item bodyček,
	\item napřiměřená hrubost
	\item zdržování hry
	\item nesportovní chování
\end{itemize}


Tresty:
\begin{itemize}
	\item \textbf{Menší faul} – nepřímý volný úder pro soupeře.
	\item \textbf{Závažný faul} – trestné střílení.
	\item \textbf{Hrubé nebo opakované porušení} – vyloučení hráče ze zápasu nebo turnaje.
\end{itemize}

Rozhodčí může posoudit situaci individuálně dle závažnosti.

\section*{8. Rozhodčí}
\begin{itemize}
	\item Každý zápas řídí rozhodčí určený organizátorem.
	\item Rozhodčí má \textbf{konečné slovo} ve všech situacích.
	\item Rozhodčí dohlíží na:
	      \begin{itemize}
		      \item dodržování pravidel,
		      \item plynulost hry,
		      \item bezpečnost hráčů.
	      \end{itemize}
	\item Nesportovní chování vůči rozhodčímu může vést k vyloučení z turnaje.
\end{itemize}

\section*{9. Výstroj}
\begin{itemize}
	\item \textbf{Rukavice jsou povinné}.
	\item Doporučena je další ochranná výstroj.
	\item Každý hráč hraje na vlastní odpovědnost.
\end{itemize}

\section*{10. Fair play}
\begin{itemize}
	\item Respektujte soupeře, spoluhráče i rozhodčí.
	\item Turnaj je koncipován jako \textbf{good feeling day}.
	\item Organizátor si vyhrazuje právo vyloučit tým nebo hráče při hrubém porušení nepsaných pravidel či principů fair play.
\end{itemize}

\section*{11. Závěrečná ustanovení}
\begin{itemize}
	\item Týmy musí být připraveny nastoupit včas.
	\item Organizátor si vyhrazuje právo upravit harmonogram.
	\item Účastí v turnaji souhlasíte s těmito pravidly.
\end{itemize}

\section*{12. Organizátorský tým}
\begin{itemize}
	\item Michal Dvořák je hlavním organizátorem.
	\item Členy organizátorského týmu jsou Kuba Dvořák, Barča Žáčková, Walter Groer.
\end{itemize}

\vfill
\begin{center}
	\textit{Hraj fér. Hraj naplno. Užij si den. Dej si pivo.}
\end{center}

\end{document}


